\chapter{Research Examples: How the Theory Reappears in Haseman Problems and Singular Integral Operators}

\section{From the boundary condition to an operator}

In Karlovich \cite[pp.~463--464, eqs.~(1.1)--(1.2)]{Karlovich2017}, $\Gamma$ is an
open star of logarithmic spirals and $1<p<\infty$. Take $G\in\SO(\Gamma)$,
$g\in L^p(\Gamma)$, and a diffeomorphism $\alpha:\Gamma\to\Gamma$ that preserves
the orientation of each branch and satisfies the source's assumptions. The Haseman
condition is
\[
 \Phi^+(\alpha(t))=G(t)\Phi^-(t)+g(t),\qquad t\in\Gamma.
\]
Here $S_\Gamma\in\mathcal B(L^p(\Gamma))$ is the Cauchy singular integral operator,
$V_\alpha f=f\circ\alpha$, and $M_Gf=Gf$. If $\varphi$ is the Cauchy density, the
Sokhotski--Plemelj formulas give
\[
 \Phi^+=P_+\varphi,\qquad \Phi^-=-P_-\varphi,\qquad
 P_\pm=\tfrac12(I\pm S_\Gamma).
\]
The sign of the negative trace explains the plus sign in the operator equation:
\[
 \bigl(V_\alpha P_++M_GP_-\bigr)\varphi=g,\qquad
 T:=V_\alpha P_++M_GP_-\in\mathcal B(L^p(\Gamma)).
\]
The course supplies the language for asking whether $T$ is invertible and for
studying its spectrum. The boundary reduction and Fredholm criterion are later
results that depend on the article's geometric and symbol assumptions.

\begin{cthejemplo}{(The identity as a case of the Haseman equation)}
On $X=L^p(\Gamma)$ consider the formula $T=V_\alpha P_++M_GP_-$ obtained above.
We check which operator results when $\alpha=\mathrm{id}_\Gamma$ and $G=1$.
Then $V_\alpha=I_X$ and $M_G=I_X$, while the definition of $P_\pm$ gives
\[
 P_++P_- =\tfrac12(I_X+S_\Gamma)+\tfrac12(I_X-S_\Gamma)=I_X.
\]
Thus $T=P_++P_-=I_X$. The Plemelj reduction here does not require
$S_\Gamma^2=I_X$.
\end{cthejemplo}

\begin{center}
\begin{tikzpicture}[node distance=5mm and 8mm]
\node[theorybox,text width=3.35cm] (b) {Haseman boundary\\$\Phi^+\circ\alpha=G\Phi^-+g$};
\node[theorybox,text width=3.35cm,right=of b] (p) {Sokhotski--Plemelj\\$\Phi^+=P_+\varphi$, $\Phi^-=-P_-\varphi$};
\node[researchbox,text width=3.35cm,right=of p] (t) {$T=V_\alpha P_++M_GP_-$\\$T\varphi=g$};
\draw[arr] (b)--(p);
\draw[arr] (p)--(t);
\end{tikzpicture}
\end{center}

\section{From the compact ideal to the Fredholm property}

Let $X$ be an infinite-dimensional complex Banach space. The closed two-sided ideal
$\mathcal K(X)\triangleleft\mathcal B(X)$ yields the Calkin algebra
$\mathcal Q(X)=\mathcal B(X)/\mathcal K(X)$. For $A\in\mathcal B(X)$,
\[
 A\text{ Fredholm}
 \quad\Longleftrightarrow\quad
 A+\mathcal K(X)\in\GL(\mathcal Q(X)).
\]
Karlovich states this equivalence in the introduction to his 2025 article
\cite[p.~2]{Karlovich2025}. Problem 1.4 explains why a quotient by a closed ideal
remains a Banach algebra; the Fredholm equivalence adds a theorem specific to operators.

\begin{cthejemplo}{(Invertible modulo compact operators without an exact inverse)}
Atkinson's criterion distinguishes an exact inverse from an inverse in the Calkin algebra.
To see this on $X=\ell^2$, let $P$ be the rank-one orthogonal projection onto
$\mathbb Ce_1$ and set $A=I_X-P$. Then
$\ker A=\operatorname{Ran}P\ne\{0\}$, so $A$ is not invertible. However,
$P\in\mathcal K(X)$ and $A+\mathcal K(X)=I_X+\mathcal K(X)$: its class is the
identity of $\mathcal Q(X)$. Moreover,
$\operatorname{Ran}(A)=(\mathbb Ce_1)^\perp$ is closed, while the kernel and
cokernel each have dimension one. Thus $A$ is Fredholm.
\end{cthejemplo}

If the \emph{Fredholm essential spectrum} is defined by
\[
 \sigma_{\mathrm{ess}}(A)
 :=\sigma_{\mathcal Q(X)}\bigl(A+\mathcal K(X)\bigr),
\]
then $0\notin\sigma_{\mathrm{ess}}(A)$ is equivalent to $A$ being Fredholm.
This does not say that $0\notin\sigma_{\mathcal B(X)}(A)$; the example $I_X-P$
shows the difference.

\begin{center}
\begin{tikzpicture}[node distance=5mm and 7mm]
\node[theorybox,text width=2.65cm] (ideal) {$\mathcal K(X)\triangleleft\mathcal B(X)$\\closed ideal};
\node[theorybox,text width=2.65cm,right=of ideal] (quot) {$\mathcal Q(X)$\\quotient};
\node[theorybox,text width=2.65cm,right=of quot] (class) {$A+\mathcal K(X)$\\invertible class};
\node[researchbox,text width=2.65cm,right=of class] (fred) {$A$ Fredholm\\Calkin criterion};
\draw[arr] (ideal)--(quot);
\draw[arr] (quot)--(class);
\draw[arr] (class)--(fred);
\end{tikzpicture}
\end{center}

On a Hilbert space $H$, $\mathcal B(H)$ and its quotient by $\mathcal K(H)$ are
$C^*$-algebras. For a general Banach space, the preceding argument needs only
the Banach algebra structure. Symbol calculations and localization seek to turn
invertibility of the class into verifiable conditions for a \emph{specific operator
algebra}; simply writing down the quotient does not produce a criterion.

\Needspace{17\baselineskip}
\section{Cauchy projections: the model case and the Haseman case}

The comparison identifies which assumption turns the formula $(I\pm S_\Gamma)/2$
into a projection and which part of the formula is used in the Haseman equation.

In the model $\Gamma=\mathbb R$ or $\mathbb T$ on $L^2(\Gamma)$, let $S_\Gamma$
be the Cauchy singular integral operator. Bastos--Fernandes--Karlovich recall
$S_\Gamma^2=I$ and that
\[
 P_\Gamma^\pm=\tfrac12(I\pm S_\Gamma)
\]
are orthogonal projections \cite[§2, p.~22]{BFK2006}. The identity
$S_\Gamma^2=I$ gives $P_\Gamma^\pm{}^2=P_\Gamma^\pm$ and
$P_\Gamma^+P_\Gamma^-=0$; self-adjointness of $S_\Gamma$ also gives
$(P_\Gamma^\pm)^*=P_\Gamma^\pm$. This is the literal connection with
projections $p=p^*=p^2$ in a $C^*$-algebra.

\begin{center}
\begin{tikzpicture}[node distance=5mm and 7mm]
\node[theorybox,text width=2.6cm] (s) {$S_\Gamma^2=I$\\$S_\Gamma^*=S_\Gamma$};
\node[theorybox,text width=2.6cm,right=of s] (def) {$P_\Gamma^\pm=(I\pm S_\Gamma)/2$};
\node[theorybox,text width=2.6cm,right=of def] (proj) {$P_\Gamma^\pm{}^2=P_\Gamma^\pm$\\$(P_\Gamma^\pm)^*=P_\Gamma^\pm$};
\node[researchbox,text width=2.6cm,right=of proj] (split) {$L^2(\Gamma)=\operatorname{Ran}P_\Gamma^+\oplus\operatorname{Ran}P_\Gamma^-$};
\draw[arr] (s)--(def);
\draw[arr] (def)--(proj);
\draw[arr] (proj)--(split);
\end{tikzpicture}
\end{center}

On the open star of Karlovich 2017, the same \emph{formula}
$P_\pm=(I\pm S_\Gamma)/2$ appears in Sokhotski--Plemelj and in
$T=V_\alpha P_++M_GP_-$, but idempotence from the preceding model cannot
be imported. If $X$ is a Banach space, $S\in\mathcal B(X)$, and
$P_\pm=(I_X\pm S)/2$, then
\[
 P_\pm^2-P_\pm=\tfrac14(S^2-I_X),\qquad
 P_+P_-=\tfrac14(I_X-S^2),\qquad P_++P_-=I_X.
\]
Indeed, the Mellin symbol of $S_\Gamma$ in the one-branch case with $p=2$
and geometric parameter $\delta=0$ is $\varepsilon_1\tanh(\pi x)$, where
$\varepsilon_1\in\{-1,1\}$ records orientation; its square is not
identically one \cite[Theorem~5.3, p.~476]{Karlovich2017}.
Thus, in this case $P_\pm$ are not algebraic projections.
The correct connection to Haseman is the \emph{decomposition of traces}
by Plemelj, not an orthogonal decomposition of the space.

\begin{cthnota}{(Where the assumptions enter)}
The identity $S_\Gamma^2=I$ is exactly what gives idempotence;
$S_\Gamma^*=S_\Gamma$ is needed for orthogonality. The model
$L^2(\mathbb R)$ or $L^2(\mathbb T)$ satisfies both properties in the cited
source. The 2017 Haseman problem acts on $L^p(\Gamma)$ over an open star and
uses only the Plemelj formula here.
\end{cthnota}

\Needspace{25\baselineskip}
\section{Unitary shifts and composition shifts}

We compare a unitary rotation with the composition operator that appears in
Haseman. The underlying space and measure determine which property can be asserted.

In \cite[§1, p.~20]{BFK2006}, the circle $\mathbb T$ acts by a unitary representation;
for $z,t\in\mathbb T$ and $\varphi\in L^2(\mathbb T)$ one writes
\[
 U:\mathbb T\to\mathcal B(L^2(\mathbb T)),
 \qquad
 (U_z\varphi)(t)=\varphi(zt).
\]
Rotation preserves Haar measure $m$: substituting $u=zt$ gives
\[
 \|U_z\varphi\|_2^2=\int_{\mathbb T}|\varphi(zt)|^2\,dm(t)
 =\int_{\mathbb T}|\varphi(u)|^2\,dm(u)=\|\varphi\|_2^2.
\]
Since $U_z^{-1}=U_{z^{-1}}$, $U_z$ is unitary. The Haseman problem instead
involves $V_\alpha f=f\circ\alpha$ on $L^p(\Gamma)$; Karlovich imposes
regularity and derivative-control assumptions on the shift on each branch
\cite[§2, p.~467]{Karlovich2017}. These conditions allow one to work with a
bounded operator, but do not assert that it is unitary.

\begin{center}
\begin{tikzpicture}[node distance=10mm]
\node[theorybox] (shift) {Common idea\\composition operator / shift};
\node[researchbox,above right=of shift] (unit) {$U_z$ on $L^2(\mathbb T)$\\unitary};
\node[researchbox,below right=of shift] (hasa) {$V_\alpha$ on $L^p(\Gamma)$\\bounded under shift assumptions};
\node[theorybox,right=22mm of unit] (uspec) {Unitary $\Rightarrow\spec(U_z)\subset\mathbb T$};
\node[theorybox,right=22mm of hasa] (fred) {Combined with $P_\pm$ and $M_G$ in the Haseman operator};
\draw[arr] (shift)--(unit);
\draw[arr] (shift)--(hasa);
\draw[arr] (unit)--(uspec);
\draw[arr] (hasa)--(fred);
\end{tikzpicture}
\end{center}

\section{Slowly oscillating functions and fibers of characters}

For $a\in C_b(\mathbb R_+)$, write
$\operatorname{osc}(a,[r,2r])=\sup_{s,t\in[r,2r]}|a(s)-a(t)|$.
Karlovich defines $\SO(\mathbb R_+)$ by requiring this oscillation to
tend to zero both as $r\to0$ and as $r\to\infty$
\cite[§2, pp.~465--466]{Karlovich2017}. For $I=[r,2r]$ and functions
$a,b\in C_b(\mathbb R_+)$,
\[
\begin{aligned}
 \operatorname{osc}(a+b,I)&\le\operatorname{osc}(a,I)+\operatorname{osc}(b,I),\\
 \operatorname{osc}(ab,I)&\le\|a\|_\infty\operatorname{osc}(b,I)
                         +\|b\|_\infty\operatorname{osc}(a,I),\\
 \operatorname{osc}(\overline a,I)&=\operatorname{osc}(a,I).
\end{aligned}
\]
If $a_n\to a$ uniformly, then
$\operatorname{osc}(a,I)\le2\|a-a_n\|_\infty+
\operatorname{osc}(a_n,I)$; first fix a large $n$ and then let $r$
approach the endpoint. Thus the class is a closed subalgebra, contains
the constants, and is stable under pointwise involution. The identity
$\|\overline a a\|_\infty=\|a\|_\infty^2$ shows that it is a unital
commutative $C^*$-algebra.

\begin{cthejemplo}{(Slow oscillation without an ordinary limit)}
We check that slow oscillation at the endpoints need not imply an
ordinary limit. On $\mathbb R_+$ consider the function
\[
 a(r)=\sin\!\bigl(\log(1+|\log r|)\bigr),\qquad r>0.
\]
It belongs to $\SO(\mathbb R_+)$. If $s,t\in[r,2r]$, then
$|\log s-\log t|\le\log2$. Put $h(u)=\log(1+|u|)$.
On the interval $I_r=[\log r,\log r+\log2]$, when it is away from $0$,
the mean value theorem gives
\[
 |h(\log s)-h(\log t)|
 \le \frac{\log2}{1+\min_{u\in I_r}|u|}.
\]
The denominator tends to infinity both as $r\to0$ and as $r\to\infty$.
Since $|\sin v-\sin w|\le|v-w|$, the oscillation of $a$ on $[r,2r]$
tends to zero at both endpoints. To see that there is no ordinary
limit, define
\[
 r_n^+=\exp\!\bigl(e^{\pi/2+2\pi n}-1\bigr),\qquad
 r_n^-=\exp\!\bigl(e^{3\pi/2+2\pi n}-1\bigr).
\]
Both sequences tend to infinity, while $a(r_n^+)=1$ and $a(r_n^-)=-1$.
Their reciprocals tend to zero and give the same two values, because
$a(r)=a(1/r)$. This function shows why boundary characters of
$\SO(\mathbb R_+)$ cannot simply be
replaced by evaluations $a(0)$ and $a(\infty)$.
\end{cthejemplo}

The maximal ideal space $\Omega(\SO(\mathbb R_+))$ contains evaluations
at $r>0$ and fibers $M_0,M_\infty$ over the endpoints
\cite[§2, p.~466]{Karlovich2017}. The Gelfand transform
$a\mapsto\widehat a$, $\widehat a(\xi)=\xi(a)$, describes values of these
characters without pretending that $a(0)$ and $a(\infty)$ exist. In
symbol calculations, these fibers are part of the parameter space;
the Gelfand map of the coefficients and the full symbol of an operator
are different maps.

\section{Generators, quotients, and symbol maps}

Karlovich \cite[eq.~(1.4), p.~3]{Karlovich2025} fixes a star-shaped
curve with cusps $\Gamma$ from (1.1)--(1.2), an admissible power weight
$w$ from §2.1, with $1<p<\infty$ and exponent
$-1/p<\alpha_{\mathrm{wt}}<1-1/p$, and considers the Banach algebra
generated by
\[
 \mathfrak B=
 \operatorname{alg}_{\mathbb R}
 \{M_a,S_\Gamma,C:\ a\in\SO^\diamond(\Gamma)\}
 \subset\mathcal B(\mathbb RL^p(\Gamma,w)).
\]
Here $M_a$ multiplies by $a$, $S_\Gamma$ is the Cauchy singular operator,
and $Cf=\overline f$ is complex conjugation. The notation
$\mathbb RL^p(\Gamma,w)$ means that the $L^p$ space is considered over
the real scalars: for $f\in L^p(\Gamma,w)$,
$C(if)=\overline{if}=-i\overline f=-iCf$, whereas
$C(rf)=rCf$ for $r\in\mathbb R$. Thus $C$ is linear over $\mathbb R$
but not over $\mathbb C$. This algebra therefore does not automatically
become a complex $C^*$-algebra, and the spectral results for complex
Banach algebras in Part 1 cannot simply be applied. The notation $\SO^\diamond(\Gamma)$
belongs to the 2025 article and does not denote the same algebra
$\SO(\mathbb R_+)$ as in the 2017 article. The 2025 article calls the
weight exponent $\alpha$; the subscript here avoids confusing it with
the shift $\alpha$ of Karlovich 2017.

\begin{center}
\begin{tikzpicture}[node distance=5mm and 8mm]
\node[researchbox,text width=3.25cm] (gen) {$M_a,S_\Gamma$ complex-linear\\$C$ real-linear};
\node[theorybox,text width=3.25cm,right=of gen] (alg) {$\mathfrak B$\\real algebra;\\norm closed};
\node[researchbox,text width=3.25cm,right=of alg] (sym) {Quotient and symbols\\criterion in a fixed algebra};
\draw[arr] (gen)--(alg);
\draw[arr] (alg)--(sym);
\end{tikzpicture}
\end{center}

Problem 1.3 teaches what it means to generate and close an algebra.
Problem 1.4 teaches how to pass to a quotient. In the 2025 article,
the additional step is to construct matrix-valued symbol maps and
prove that they detect the Fredholm property. Theorem~6.2 treats the
complex subalgebra
$\mathfrak A=\operatorname{alg}_{\mathbb C}
\{M_a,S_\Gamma,CS_\Gamma C:a\in\SO^\diamond(\Gamma)\}$;
the complex matrix representation of Theorem~7.1 then allows the
criterion for $N\in\mathfrak B$ to be formulated using the symbol map
$Y$ of Theorem~7.2 \cite[pp.~25--26]{Karlovich2025}. The latter
requires membership in $\mathfrak B$ and the article's assumptions
on the curve, weight, and coefficients; it does not by itself certify
a Haseman operator with a shift.

\begin{cthnota}{(Four related but distinct maps)}
An algebra homomorphism preserves sums and products. A
$*$-homomorphism also preserves the involution. The Gelfand transform
$a\mapsto\widehat a$ of a commutative algebra evaluates at its
characters. A Fredholm symbol map is constructed for a particular
operator algebra, often after passing modulo compact operators, and
may take matrix values. These constructions share algebraic
operations, but their domains, codomains, and assumptions cannot be
interchanged.
\end{cthnota}

\Needspace{20\baselineskip}
\section{The full route from the problem to a criterion}

In the diagram, blue marks the language of the course, green the
research object, and the arrow toward symbols presupposes that the
corresponding calculus has been constructed for the particular
algebra. Here $X=L^p(\Gamma)$, so $T+\mathcal K(X)$ is a class in the
Calkin algebra associated with the same underlying space. The
criterion of Karlovich 2017 also requires the assumptions of its
Theorem 7.1 \cite[pp.~487--489]{Karlovich2017}.

\begin{center}
\begin{tikzpicture}[node distance=5mm and 7mm]
\node[researchbox,text width=3.25cm] (h) {[research]\\Haseman boundary};
\node[researchbox,text width=3.25cm,right=of h] (p) {[research]\\Sokhotski--Plemelj};
\node[researchbox,text width=3.25cm,right=of p] (t) {[research]\\$T=V_\alpha P_++M_GP_-$};
\node[theorybox,text width=3.25cm,below=of t] (b) {[course]\\$T\in\mathcal B(L^p(\Gamma))$};
\node[theorybox,text width=3.25cm,left=of b] (q) {[course + bridge]\\$T+\mathcal K(X)$ in $\mathcal Q(X)$};
\node[researchbox,text width=3.25cm,left=of q] (f) {[bridge]\\invertibility of the class\\$\Longleftrightarrow$ Fredholm};
\node[researchbox,text width=3.25cm,below=of f] (s) {[research]\\symbols and localization};
\node[researchbox,text width=3.25cm,right=of s] (c) {[research]\\Fredholm criterion\\under article assumptions};
\draw[arr] (h)--(p);
\draw[arr] (p)--(t);
\draw[arr] (t)--(b);
\draw[arr] (b)--(q);
\draw[arr] (q)--(f);
\draw[arr] (f)--(s);
\draw[arr] (s)--(c);
\end{tikzpicture}
\end{center}

The last row describes how the Fredholm property is \emph{checked} in
the literature. It does not claim that every Fredholm operator
automatically has the particular symbol of the article, or that the
2017 criterion applies to a 2025 operator with cusps.

\Needspace{19\baselineskip}
\section{A reading path for the research manuscript}

When an operator construction appears in the manuscript, it is
helpful to ask the following questions in order:
\begin{enumerate}
\item In which Banach algebra does the operator live?
\item What is the identity, and what are the relevant generators?
\item Are we working exactly or modulo an ideal, usually $\K(X)$?
\item Is the subalgebra under consideration commutative, or does it only
      have a commutative central subalgebra?
\item Do characters, a maximal ideal space, or a Gelfand transform appear?
\item Is there a complex $C^*$-algebra structure (typically on $L^2$)
      that permits the use of normal elements, unitaries, projections,
      or functional calculus?
\item Which part of the Fredholm criterion is equivalent to
      invertibility in a quotient or of a symbol?
\end{enumerate}

These questions distinguish the common language from the assumptions
that authorize each theorem. For example, a published criterion for
an operator algebra on a curve with cusps does not automatically apply
to an operator with a shift on a different geometry.
